\section{関連研究}\label{sec:related}

\paragraph{イジングマシンとハイブリッド.} CIM\cite{wang2013,inagaki2016,mcmahon2016}のほか，カオス振幅制御（CAC）\cite{leleu2019}，シミュレーテッド分岐（SB）\cite{goto2019}など，連続変数の力学で解を探す方式が提案されている．CIM の MAX-CUT 性能は SDP 緩和\cite{goemans1995}や SA と比べて評価されてきた\cite{haribara2015}．本論文はイジングマシンの出力を後段の古典的な探索に渡す使い方を対象にし，渡す候補の選び方に焦点を当てる．

\paragraph{MAX-CUT の発見的解法.} G-set\cite{helmberg2000}に対しては，ブレイクアウト局所探索（BLS）\cite{benlic2013}，メメティックアルゴリズム\cite{wu2012}，複数の探索演算子を使う MOH\cite{ma2017}などが高い性能を示している．Dunning ら\cite{dunning2018}は多数の解法を体系的に比較した．本論文の目的は最良値の更新ではなく，CIM と後段の探索をつなぐ部分の設計である．

\paragraph{初期解の選び方と地形.} 局所最適解の網（local optima network）\cite{ochoa2008}は，探索空間の地形を局所最適解とその間の遷移で表し，探索の難しさを説明する．本論文の障壁・経路長はこれと同じ問題意識を，個々の初期解について厳密に測ったものである．強化学習で解を改善する方法\cite{barrett2020}も，初期状態からどこへ到達できるかを学習している．計算予算の配分の観点では，短い試行で候補を絞る successive halving\cite{karnin2013,jamieson2016}が，本論文の「試し打ち」に相当する強い比較相手である．

\paragraph{スペクトル的な手法.} MAX-CUT と隣接行列の最小固有値・固有ベクトルの関係は，スペクトル的な丸めとして知られている\cite{trevisan2012}．CIM の増幅初期は線形写像で近似でき，結合行列の最大固有値の固有モードが先に成長するので，CIM の出力がそのモードに偏ることが考えられる．本論文の固有ベクトル指標は，この偏りを候補ごとに測るものである．
